\section{Calibration in the chart}
\label{sec:svicalib}

Chapter~\ref{ch:lqd} settled the anatomy of a slice objective --- units,
band targets, regularization (\cref{sec:calibration}).  Working in the
volatility chart changes two of the choices, and this section treats them
in turn.  The residual can now live directly in implied volatility, since
the model produces $w(k)$ without an inversion: with quotes
$(k_i,\sigma_i)$ and weights $\mu_i$, the data term is
\begin{equation}
\sum_i \mu_i\,\big(I(k_i;\theta)-\sigma_i\big)^2,
\qquad I(k;\theta)=\sqrt{w(k;\theta)/\tau},
\label{eq:sviobjective}
\end{equation}
optionally replaced by the band or haircut targets of \eqref{eq:band} and
\eqref{eq:haircut} evaluated in volatility space.  And the chart itself
becomes a decision, because the coordinates no longer guarantee validity
for free.

\subsection{The measure a fit integrates against}
\label{sec:sviweights}

Idealize \eqref{eq:sviobjective} as a continuum problem: minimize
$\int(I-\sigma)^2\,\dd\mu(k)$ for some measure $\mu$ on the strike axis.
A finite quote set turns the integral into a sum, and the weights are the
quadrature rule.  Equal weights make $\mu$ the empirical measure of the
\emph{listing} --- the exchange's strike schedule, dense at the money and
sparse in the wings.  That measure was chosen by a listings committee: a
tight cluster of near-duplicate ATM quotes outvotes an isolated wing quote
by sheer count, and the fit inherits the strike schedule as an accidental
prior.  Adding more ATM listings, which adds information, then makes the
fitted wings \emph{worse}.

A declared alternative weights strike ranges by the time value they
carry, $\dd\mu=\mathrm{TV}(k)\,\dd k$, where $\mathrm{TV}(k)$ is the
normalized price of the out-of-the-money option at $k$ --- the part of
the premium that is genuinely optionality rather than intrinsic value,
the natural gauge of how much economic content a strike carries.  Let
$C_i$ be quote $i$'s cell in the
Voronoi partition of the quoted interval --- half the gap to each
neighbor, one-sided at the ends --- with width $|C_i|$.  The scheme
\begin{equation}
\mu_i\;\propto\;\mathrm{TV}_i\,|C_i|
\label{eq:tvweights}
\end{equation}
is then a genuine discretization of the declared measure.

\begin{proposition}[Weights as a quadrature]
\label{prop:quadrature}
For any continuous integrand $f$, the sum $\sum_i\mathrm{TV}_i|C_i|f(k_i)$
is a Riemann sum of $\int\mathrm{TV}(k)f(k)\,\dd k$ over the Voronoi
partition, converging under refinement.  In particular the weighted loss
with \eqref{eq:tvweights} discretizes the continuum time-value objective;
on a uniform strike grid all $|C_i|$ coincide and the scheme reduces to
pure time-value weighting.
\end{proposition}

\begin{proof}
The cells partition the quoted interval and $k_i\in C_i$, so the sum is a
Riemann sum over that partition; refinement convergence is the standard
statement.  The uniform case is arithmetic.
\end{proof}

Two conventions keep the quadrature usable on real ladders, both stated
once.  The cell of an isolated far-wing quote can be arbitrarily wide, so
the width factor is capped (the reference implementation uses ten times
the mean spacing): beyond the cap the scheme deliberately
under-integrates the deep wing rather than let one quote carry a tenth of
the objective.  And the weights are normalized to unit mean, so switching
schemes redistributes the data term without rescaling it against the
regularization.  The chapter's comparisons use equal weights throughout
--- the point here is not that one measure is right, but that a weight
vector \emph{is} a measure, and should be chosen as one.

\subsection{Two charts}
\label{sec:svistructural}

The optimizer never handles a constrained vector.  In the historical
chart, $\theta\in\R^5$ maps to
\begin{equation}
a_S=\theta_1,\quad
b_S=\operatorname{softplus}(\theta_2),\quad
\rho_S=\tanh(\theta_3),\quad
m_S=\theta_4,\quad
s_S=e^{\theta_5},
\label{eq:rawlift}
\end{equation}
with $\operatorname{softplus}(\theta)=\log(1+e^{\theta})$, so every
finite $\theta$ is a structurally valid hyperbola --- the first row of
\cref{tab:tiers} and nothing more.  The floor can still go
negative and a wing can still exceed the cap; those are fenced by the
hinge rows \eqref{eq:svipenalties}, which the optimizer can trade against
the data term.  This is the same constraint-elimination idea as the
logistic chart of \cref{sec:charts}, applied to the weaker guarantee.

The idea extends to the conditions the hinges police.  Parameterize the
slice by the quantities the inequalities are \emph{about} --- the two wing
slopes, the vertex location, the floor, and the vertex curvature --- each
lifted so that every finite coordinate is admissible:
\begin{equation}
\beta_L=\beta_{\max}\Lambda(\theta_1),\quad
\beta_R=\beta_{\max}\Lambda(\theta_2),\quad
k^*=\theta_3,\quad
w^*=\operatorname{softplus}(\theta_4),\quad
\kappa^*=e^{\theta_5},
\label{eq:structlift}
\end{equation}
with $\Lambda$ the logistic function of \cref{sec:intro}.  The raw
parameters are recovered exactly by inverting \cref{prop:svigeom}:
\begin{equation}
b_S=\frac{\beta_L+\beta_R}{2},\quad
\rho_S=\frac{\beta_R-\beta_L}{\beta_R+\beta_L},\quad
s_S=\frac{b_S(1-\rho_S^2)^{3/2}}{\kappa^*},\quad
m_S=k^*+\frac{s_S\rho_S}{\sqrt{1-\rho_S^2}},\quad
a_S=w^*-b_Ss_S\sqrt{1-\rho_S^2}.
\label{eq:structinv}
\end{equation}
Under this \emph{structural chart} every finite iterate has
$w(k)\ge w^*>0$ and both wings strictly inside $(0,\beta_{\max})$: the
second and third rows of \cref{tab:tiers} hold by construction, and the
hinge rows \eqref{eq:svipenalties} are identically zero along the entire
optimization path.  What the chart does \emph{not} deliver is the fourth
row --- the Vogt slice of \cref{sec:svibelly} is reachable in either
chart, and the dense-grid check remains the deciding instance.  The two
charts describe the same family, and on the SPY December quotes they
reach the same slice to $\MacSviChartAgreeBp$ vol bp
(\cref{fig:svistructural}); what changes is the geometry the solver
works in, with the kinks of the hinge rows removed exactly where a
trust-region step would otherwise stall on them.

\begin{figure}[htbp]
\centering
\paperfig{0.97\textwidth}{fig_svi_structural.pdf}
\caption{The structural chart.  (a)~The wing lift
$\beta=\beta_{\max}\Lambda(\theta)$: every finite coordinate lands
strictly between $0$ and $\beta_{\max}$, itself strictly below Lee's
bound --- admissibility is a property of the chart, not a constraint on
the step.  (b)~The five structural coordinates read on the SPY December
fit: vertex $(k^*,w^*)=(\MacSviSpyKstar,\MacSviSpyWstar)$, osculating
parabola of curvature $\kappa^*=\MacSviSpyKappastar$, and the two
asymptote slopes.  The optimizer moves exactly the quantities the
inequalities constrain.}
\label{fig:svistructural}
\end{figure}

\subsection{The residual Jacobian, in layers}
\label{sec:svijacobian}

The least-squares solvers of both charts consume an analytic Jacobian,
and it assembles by the chain rule in three layers.  The raw-parameter
partials of the hyperbola follow from \eqref{eq:svi} with
\eqref{eq:svishort}:
\begin{equation}
\frac{\partial w}{\partial a_S}=1,\quad
\frac{\partial w}{\partial b_S}=\rho_Sx+r,\quad
\frac{\partial w}{\partial\rho_S}=b_Sx,\quad
\frac{\partial w}{\partial m_S}=-b_S\Big(\rho_S+\frac{x}{r}\Big),\quad
\frac{\partial w}{\partial s_S}=\frac{b_Ss_S}{r}.
\label{eq:svipartials}
\end{equation}
The data rows live in volatility, so each is scaled by
\begin{equation}
\frac{\partial I}{\partial w}=\frac{1}{2\sqrt{w\tau}}=\frac{1}{2\tau I}.
\label{eq:ivchain}
\end{equation}
Finally the lift: for \eqref{eq:rawlift} the three nontrivial factors are
\begin{equation}
\frac{\dd b_S}{\dd\theta_2}=1-e^{-b_S},\qquad
\frac{\dd\rho_S}{\dd\theta_3}=1-\rho_S^2,\qquad
\frac{\dd s_S}{\dd\theta_5}=s_S,
\label{eq:liftchain}
\end{equation}
while the structural chart pushes the same raw-space rows through the
closed-form $5\times5$ matrix
$\partial(a_S,b_S,\rho_S,m_S,s_S)/\partial\theta$ obtained by
differentiating \eqref{eq:structinv} along \eqref{eq:structlift}.  Hinge
rows contribute their linear part where the violation is strictly
positive and a zero row otherwise.  Nothing here is deep; the point of
recording the layers is that the chart is a change of variables like any
other, and its cost is one matrix product per iterate rather than a
constrained solver.
